\section{Proposed Model or Transformer}
\label{sec:model}

\subsection{Backbone}

The proposed backbone is \textbf{DINOv3 ViT-L/16} \cite{simeoni2025dinov3}, loaded from the Hugging Face hub as \texttt{\backboneId}. It is a vision transformer \cite{dosovitskiy2021vit} with 24 layers, 1,024-dimensional tokens, $16\times16$ pixel patches and four register tokens \cite{darcet2024registers}, about 300 million parameters, trained without labels by self-distillation on a curated set of about 1.7 billion images and distilled from a 7-billion-parameter teacher. Its role in the study is that of a feature extractor that is fine-tuned end to end: the image with the highlighted query enters as a sequence of patch tokens, and the relation decision is read from the output tokens.

Four properties make it suitable, each tied to a property of the data. Global self-attention relates every patch to every other patch in every layer, whereas a convolutional network such as ResNet \cite{he2016resnet} connects distant regions only in its deepest layers; this matters because a relation is a property of primitives that may lie far apart in the image. Self-supervised training on images alone yields dense, stroke-level features, and DINOv3's \emph{Gram anchoring} is designed to keep the patch-level features from degrading over long training \cite{simeoni2025dinov3}; caption-aligned encoders such as CLIP \cite{radford2021clip} learn what captions describe, and captions do not describe angles. The model accepts inputs of varying size, so diagrams are processed at a resolution that keeps thin lines and small marks without distortion. And its features are strong when frozen \cite{oquab2024dinov2,simeoni2025dinov3}, which allows a frozen-feature baseline on the small real training set (\nFiguresTrain{} figures) as well as full fine-tuning.

\textbf{DINOv2 ViT-L/14} \cite{oquab2024dinov2} (\texttt{\fallbackId}) is the named fallback and comparison: it is peer-reviewed, whereas DINOv3 is at present a preprint, and it is run under the identical recipe in one comparison (Section~\ref{sec:baselines}). At the same input size its 14-pixel patches yield 1,024 patch tokens against 784 for DINOv3, a difference the comparison must keep in mind. The backbone identifier is a single configuration value, so switching between them changes nothing else.

\subsection{Feature pooling and head}

From the backbone's output the model drops the register tokens and forms a pooled feature by concatenating the CLS token with the mean of the patch tokens (2,048 dimensions), so that both a global summary and the average local evidence reach the head. The head is a LayerNorm followed by one linear layer with six outputs and a sigmoid per output: a multi-label classifier, because a segment pair can be both perpendicular and equal in length. The loss is binary cross-entropy over the applicable relations only (the mask zeroes the others), with the per-relation positive weights of Section~\ref{sec:preprocessing}. Gradient checkpointing \cite{chen2016checkpointing} is enabled so that the full fine-tune fits in a few gigabytes of GPU memory at \inputLongSide{} px. Every choice named in this section is one entry of the configuration file summarised in Table~\ref{tab:config}.

\subsection{Query-as-tokens variant}

To test whether the highlight overlay, rather than any other part of the design, is what makes the task learnable, one comparison model receives the query not as a highlight but as tokens: the query kind and the point names, hashed to token indices, are embedded and projected to a single query vector that cross-attends over the CLS and patch tokens, and the attended vector is concatenated with the pooled feature before the head. This model can only answer by finding the named letters in the image. It is trained with the same budget but without the probe stage, which it cannot use because its query encoder has no frozen features to probe.
